\chapter{The Constant the Apparatus Can Touch}

The construction has closed. What follows is not another summit inside it, but its use. The reader now has an
instrument with a count, a charge, a mass, a sign, a gauge, a trace, a cost ledger, and a discipline for saying
which of those things is built, proved, read, reserved, or handed in from outside. To measure the fine structure
constant with that instrument is to ask a very particular question of the world: how large is the electromagnetic
coupling when every unit has been stripped away?

The answer is a pure number. In SI language it is written
\[
  \alpha = \frac{e^2}{4\pi\varepsilon_0\hbar c},
\]
and its inverse is close to 137. The current CODATA 2022 recommended value is
\[
  \alpha^{-1}=137.035\,999\,177(21),
\]
or, equivalently, \(\alpha \approx 7.297\,352\,5643\times 10^{-3}\). That number is not an input to the
measurement. It is the benchmark the finished measurement will face. The apparatus must not tune itself toward
137 because the number is famous, or toward CODATA because the table is authoritative. Authority is not a rung.
The measurement has to produce its own count, carry its own corrections, and then stand beside the recommended
value with its uncertainty visible.

\section{What Alpha Measures}

The constant is dimensionless, and that is the first mercy. A dimensional constant always carries the hand's
choice of units in its clothing. A meter, a second, a coulomb, a joule: each is a gauge hung on the world. A
dimensionless constant survives the stripping of those gauges. Change the units and the number stays put. That is
why \(\alpha\) is a proper target for this apparatus. It is not asking how many of our centimeters fit a path, or
how many of our seconds fit a tick. It is asking for the ratio that remains after the gauges have cancelled.

Read in the language already built, \(\alpha\) compares three things. It compares the electron's charge, the
speed at which electromagnetic influence propagates, and the quantum of action that says how much phase is paid
for a turn. Charge is the loop count. The quantum of action is the cost of phase, the price of one accountable
turn. The speed of light is the limiting reach of the physical reading, the conversion between a spatial gap and
a temporal one. Put them together and the constant says, in one number, how strongly the counted charge couples
to the field whose signals travel at \(c\), measured against the phase-cost \(\hbar\). The words are physical;
the number they aim at has no unit.

This matters because the apparatus must not measure \(\alpha\) by measuring \(e\), \(\hbar\), and \(c\) as if
they were three ordinary lengths on a table. Since the SI redefinition, \(e\), \(h\), and \(c\) are defining
constants, fixed exactly by convention. If the apparatus only plugs exact constants into the formula above, it
has measured nothing. It has merely unfolded the unit system. To touch \(\alpha\), the instrument needs a route
where a real physical reading enters: a recoil, a magnetic moment, a Hall resistance, a spectral interval, some
place where the world answers and the machine has to carry the answer with its uncertainty.

\section{What Lean Does in the Run}

Lean is the checked notebook. It is a proof assistant and a small programming language whose central habit is
simple: every object has a type, every field has a name, every theorem has to typecheck, and every computation
has to be made from definitions the compiler can see. It is not the laser, not the atom cloud, and not the
constants adjustment. It does not make a bad measurement good. What it does is prevent the arithmetic and the
ledger from quietly changing shape after the data arrive.

For this use, the reader should treat Lean as four things at once. It is a label maker: measured quantities,
given constants, readings, and proved identities are not allowed to masquerade as one another. It is a calculator:
the recoil ratio, correction sum, bridge product, square root, reciprocal, and uncertainty propagation can be
computed in one place. It is a receipt: a theorem such as \(\texttt{rfl}\) says that the field in the final report
really is the definition it claims to be. And it is a guardrail: if the structure of the report is changed, Lean
forces every dependent definition to face the change.

The live apparatus may index a trace rung by its own constructed number type, not by Lean's built-in natural
numbers:

\begin{quote}\small
\begin{verbatim}
class TRACED
    (Value : Type i)
    (Carrier : CarrierProcess Value)
    (rung : Measurement.Natural) where
  tape : CompilerTape
\end{verbatim}
\end{quote}

That is the strict reading of a rung: the number is not borrowed from civilization but built inside the
measurement vocabulary. The alpha computation has a different surface. Its lab readings are real-valued
quantities with uncertainties, so the manual uses floating-point numerals for the closing arithmetic and carries
the apparatus discipline in the record types. The important separation is this: Lean's \(\texttt{Float}\) is the
calculator's decimal hand, while the surrounding structures say what kind of thing each decimal is.

The companion file is \(\texttt{device/Measurement/AlphaManual.lean}\). It begins by naming the grades and the
data cells:

\begin{quote}\small
\begin{verbatim}
inductive Grade where
  | measured
  | given
  | reading
  | proved

structure Datum where
  value : Float
  sigma : Float
  grade : Grade
\end{verbatim}
\end{quote}

This is the first instruction for computation. Do not enter a number naked. Enter a value, a standard
uncertainty, and a grade. A measured recoil, a given Rydberg constant, and a final reading of
\(\alpha^{-1}\) may all look like decimals, but they are not the same kind of record.

\section{The Recoil Route}

The cleanest route for the apparatus is recoil. Choose an atom \(X\), usually rubidium or cesium, whose relative
atomic mass is known with a separate uncertainty. Kick it with light of known wavelength, make it absorb and emit
photons in a controlled sequence, and measure the tiny change in its velocity. A photon has momentum \(h/\lambda\).
When an atom takes that momentum, the recoil velocity contains the ratio \(h/m_X\). The experiment therefore does
not measure \(h\) alone or \(m_X\) alone. It measures their quotient, a dimensionful ratio whose units are later
cancelled by other already audited constants.

The bridge from recoil to \(\alpha\) is
\[
  \alpha^2 =
  \frac{2R_\infty}{c}\,
  \frac{A_{\rm r}(X)}{A_{\rm r}(e)}\,
  \frac{h}{m_X}.
\]
Here \(R_\infty\) is the Rydberg constant, \(A_{\rm r}(X)\) is the relative atomic mass of the atom used in the
recoil measurement, and \(A_{\rm r}(e)\) is the relative atomic mass of the electron. The apparatus measures the
last factor, \(h/m_X\), and imports the other factors as named external inputs, each with its own uncertainty.
Then it takes the square root. That square root is not decoration. It means every fractional uncertainty in the
product arrives in \(\alpha\) at half strength, because the measured product is \(\alpha^2\).

The route is honest because it has a visible seam. The world supplies the recoil. The laboratory supplies the
gauge for light, timing, geometry, and atom choice. The constants ledger supplies \(R_\infty\) and the mass
ratios. The apparatus combines them only after each has been named for what it is. The final number is not pure
construction and not pure authority. It is a measured recoil folded through a carefully labelled dictionary into
the dimensionless coupling.

\section{The Grade of the Target}

The grade is therefore mixed. The formula relating \(\alpha^2\) to \(h/m_X\) is not a new empirical guess; it is
the accepted physical bridge between atomic recoil, the Rydberg scale, and electromagnetic coupling. The recoil
reading is MEASURED. The imported constants are GIVEN, not by the machine's own run but by outside metrology. The
identification of the resulting dimensionless number with the fine structure constant is a READING laid over
that bridge, a reading so standard that physics treats it as ordinary, but still a reading in the apparatus'
strict ledger.

That grade is the first instruction to the user. Do not begin by asking the apparatus to prove why the universe
chose \(1/137.035999...\). It will not. Begin by asking it to produce an audited value of the coupling: a number,
an uncertainty, a correction ledger, and a comparison. The apparatus can do that. It can make \(\alpha\) pass
through its own discipline: repeatability, gauge, noise, trace, residue, tange, funge, and final handoff to the
reader who decides whether the account is good enough.

The Lean rule is the same rule in a harder font. A theorem may prove that the corrected recoil value equals the
raw value plus the correction sum. It may prove that the report's reciprocal-alpha field is exactly the reciprocal
computed from the bridge. It does not prove that the wavefront correction was physically right, or that the atom
cloud was cold enough, or that the blind was honestly held. Those are laboratory facts. Lean keeps the route from
being rewritten while those facts are being carried.
